\documentclass[../DoAn.tex]{subfiles}
\begin{document}

\begin{center}
    \Large{\textbf{TÓM TẮT NỘI DUNG ĐỒ ÁN}}\\
\end{center}
\vspace{1cm}
Nhận diện chữ số viết tay là một bài toán quan trọng trong lĩnh vực thị giác máy tính với nhiều ứng dụng thực tiễn như xử lý séc ngân hàng tự động, phân loại thư từ tại bưu điện, và số hóa tài liệu giấy. Mặc dù đã được nghiên cứu từ lâu, bài toán này vẫn đặt ra thách thức do sự đa dạng trong phong cách viết tay, độ nghiêng, và chất lượng ảnh. Hiện nay có hai hướng tiếp cận chính để giải quyết vấn đề này: học máy truyền thống, và học sâu. \\
\indent Trong đồ án này, em tìm hiểu về hai phương pháp: Support vector machine (SVM), đại diện cho học máy truyền thống; và Convolutional neural network (CNN), đại diện cho học sâu; cùng ứng dụng của chúng trong bài toán phân loại chữ số viết tay trong dataset MNIST. Em chọn hướng tiếp cận này để hiểu về cả học máy truyền thống và học sâu, đồng thời nhận biết sự khác biệt trong cách mỗi phương pháp được ứng dụng trong thực tế. Sau khi tìm hiểu cơ sở lý thuyết của 2 phương pháp, em bắt đầu quá trình huấn luyện 2 model tương ứng với mỗi phương pháp. Quá trình huấn luyện bao gồm tải và chia tập dữ liệu thành 2 phần: training và test, tiền xử lý dữ liệu cho mỗi mô hình, xây dựng mô hình, huấn luyện và đánh giá. Kết quả đánh giá cho thấy mô hình CNN có hiệu năng cao hơn, với độ chính xác 98.99\% chỉ sau 51.28 giây, so với mô hình SVM, với độ chính xác 97.92\% sau 194.61 giây. 


\begin{flushright}
Sinh viên thực hiện\\
\begin{tabular}{@{}c@{}}
\textit{(Ký và ghi rõ họ tên)}
\end{tabular}
\end{flushright}

\end{document}